\documentclass[12pt,a4paper]{article}
\usepackage[utf8]{inputenc}
\usepackage[spanish,es-tabla]{babel}
\usepackage{geometry}
\usepackage{graphicx}
\usepackage{multirow}
\usepackage{tabularx}
\usepackage{mathptmx}
\usepackage{fancyhdr}
\usepackage{array}
\usepackage[table]{xcolor}
\usepackage{titlesec}
\usepackage{setspace}
\usepackage{microtype}
\usepackage{lastpage}
\usepackage{booktabs}
\usepackage{enumitem}
\usepackage{amsmath}
\usepackage{amsfonts}
\usepackage{amssymb}

\usepackage{float}
\usepackage{placeins}
\usepackage{tikz}
\usepackage{pgfplots}
\pgfplotsset{compat=1.17}
\usetikzlibrary{arrows.meta,positioning,shapes.geometric,calc}

\graphicspath{{images/}{../images/}{./}}

\geometry{
    top=4.7cm,
    bottom=2.5cm,
    left=2.5cm,
    right=2.5cm,
    headheight=3.9cm,
    headsep=0.4cm
}

\setstretch{1.15}
\setlength{\parindent}{0pt}
\setlength{\parskip}{0.45em}
\setlength{\tabcolsep}{4pt}

\titleformat{\section}
  {\normalfont\fontsize{12}{14.5}\bfseries}{\thesection.}{0.6em}{\uppercase}
\titlespacing*{\section}
  {0pt}{1.8ex plus 0.4ex minus 0.2ex}{0.6ex plus 0.2ex}

\titleformat{\subsection}
  {\normalfont\fontsize{11.5}{13.5}\bfseries}{\thesubsection.}{0.5em}{}
\titlespacing*{\subsection}
  {0pt}{1.4ex plus 0.3ex minus 0.1ex}{0.4ex plus 0.1ex}

\titleformat{\subsubsection}
  {\normalfont\fontsize{10.5}{12.5}\bfseries\itshape}{\thesubsubsection.}{0.4em}{}
\titlespacing*{\subsubsection}
  {0pt}{1.1ex plus 0.2ex minus 0.1ex}{0.3ex plus 0.1ex}

\newcolumntype{Y}{>{\centering\arraybackslash}X}
\newcolumntype{L}{>{\raggedright\arraybackslash}X}
\newcolumntype{C}{>{\centering\arraybackslash}X}

\definecolor{headerblue}{RGB}{26, 54, 93}
\definecolor{darkblue}{RGB}{43, 108, 176}
\definecolor{lightgraybg}{RGB}{248, 248, 248}
\definecolor{tableheader}{RGB}{232, 238, 246}
\definecolor{navyheader}{RGB}{0, 32, 96}
\definecolor{lineblue}{RGB}{0, 102, 204}
\definecolor{areablue}{RGB}{220, 236, 252}
\definecolor{lcsred}{RGB}{217, 56, 41}
\definecolor{lcigreen}{RGB}{46, 125, 50}
\definecolor{lciorange}{RGB}{230, 81, 0}
\definecolor{ishikawared}{RGB}{160, 0, 0}

\tikzset{
    catbox/.style={draw, fill=tableheader, font=\footnotesize\bfseries, minimum width=2.8cm, minimum height=0.65cm, align=center},
    efectobox/.style={draw=ishikawared!90!black, line width=1.2pt, fill=ishikawared, text=white, font=\small\bfseries, align=center, inner sep=6pt, text width=3.8cm},
    causetext/.style={font=\scriptsize, align=right, text=black!90}
}

\newcommand{\docCodigo}{TI-SIGE-ENT-F3-015}
\newcommand{\docVersion}{1.0}
\newcommand{\docProceso}{IMPLEMENTACIÓN DE SOFTWARE / EXPORTACIÓN DOCUMENTAL SGC}
\newcommand{\docElaboracion}{09/10/2026}
\newcommand{\docRevision}{09/10/2026}
\newcommand{\docTitulo}{MOTOR DE EXPORTACIÓN OFICIAL EN PDF Y WORD}
\newcommand{\docOrganizacion}{FACULTAD DE INGENIERÍA EN SISTEMAS, ELECTRÓNICA E INDUSTRIAL}
\newcommand{\docSubOrganizacion}{CARRERA DE INGENIERÍA DE SOFTWARE}

\newcommand{\encabezadofrecuente}{%
    \renewcommand{\arraystretch}{1.15}%
    \noindent\makebox[\textwidth][c]{%
    \begin{tabularx}{\textwidth}{|>{\centering\arraybackslash}m{2.8cm}|Y|>{\centering\arraybackslash}m{4.8cm}|}
        \hline
        \multirow{4}{2.8cm}{\centering\raisebox{-1.15cm}{\includegraphics[width=2.5cm,height=2.0cm,keepaspectratio]{images/logo.png}}}
        & \textbf{\fontsize{9}{11}\selectfont \docOrganizacion} 
        & \textbf{\fontsize{8}{10}\selectfont CÓDIGO:} \fontsize{8}{10}\selectfont \docCodigo \\
        \cline{3-3}
        & \textbf{\fontsize{8.5}{10.5}\selectfont \docSubOrganizacion} 
        & \textbf{\fontsize{8}{10}\selectfont VERSIÓN:} \fontsize{8}{10}\selectfont \docVersion \\
        \cline{2-3}
        & \textbf{\fontsize{8}{10}\selectfont PROCESO:} \fontsize{8}{10}\selectfont \docProceso 
        & \textbf{\fontsize{8}{10}\selectfont ELABORACIÓN:} \fontsize{8}{10}\selectfont \docElaboracion \\
        \cline{2-3}
        & \textbf{\fontsize{8.5}{10.5}\selectfont \docTitulo} 
        & \textbf{\fontsize{8}{10}\selectfont ÚLTIMA REVISIÓN:} \fontsize{8}{10}\selectfont \docRevision \\
        \hline
    \end{tabularx}%
    }%
}

\pagestyle{fancy}
\fancyhf{}
\fancyhead[C]{\encabezadofrecuente}
\fancyfoot[C]{\fontsize{10}{12}\selectfont Página \thepage\ de \pageref{LastPage}}
\renewcommand{\headrulewidth}{0pt}

\begin{document}

\noindent{\fontsize{10.5}{12.5}\selectfont \textbf{Integrantes:} Heidi Villavicencio \quad$\vert$\quad Steeven Toala \quad$\vert$\quad Nixon Hurtado \quad$\vert$\quad Andrea Vásquez}

\vspace{0.25cm}

\section{Título}
\textbf{Denominación Oficial del Entregable:} Motor de Exportación Oficial en PDF y Word -- Código EDT: 3.15

\section{Introducción}
Aunque la interacción diaria y el ciclo de revisión del sistema SIGE se efectúan dentro de la plataforma web, los requerimientos institucionales de la Universidad Técnica de Ambato exigen la entrega física y archivo documental de expedientes legalizados en formatos oficiales (\texttt{.pdf} vectorial no editable y \texttt{.docx} interoperable). Cualquier discrepancia de márgenes, encabezados institucionales, códigos de formato o ubicación de firmas respecto a las plantillas \texttt{UTA-SGC-A-2-1-P7-T1} y \texttt{T2} representaría una no conformidad de calidad. El presente entregable documenta el motor de generación documental de lado del servidor (\textit{Server-Side Document Export Pipeline}) que renderiza e imprime documentos idénticos al estándar institucional.

\section{Objetivo}
\subsection{Objetivo General}
Desarrollar y validar el motor de exportación de documentos oficiales en formatos PDF y Word idénticos al estándar institucional del SGC de la UTA.
\subsection{Objetivos Específicos}
\begin{itemize}[leftmargin=1.5em, itemsep=2pt]
    \item Implementar la generación de PDF vectorial de alta precisión mediante Puppeteer y estilos CSS Paged Media (@page A4).
    \item Desarrollar el generador de archivos Microsoft Word (\texttt{.docx}) estructurado mediante plantillas OpenXML.
    \item Incorporar códigos QR con URL de verificación y estampas digitales de firmas.
\end{itemize}

\section{Alcance}
\subsection{Límites del Entregable (In-Scope)}
\begin{itemize}[leftmargin=1.5em, itemsep=2pt]
    \item Exportación oficial para el Plan de Trabajo Docente (\texttt{P7-T1}) y el Informe de Cumplimiento (\texttt{P7-T2}).
    \item Preservación exacta de logotipos, códigos de control SGC, tablas de actividades y bloques de firmas.
\end{itemize}
\subsection{Exclusiones (Out-of-Scope)}
\begin{itemize}[leftmargin=1.5em, itemsep=2pt]
    \item Impresión física en papel sobre impresoras matriciales o térmicas.
\end{itemize}

\section{Definiciones}
\begin{itemize}[leftmargin=1.5em, itemsep=2pt]
    \item \textbf{Puppeteer:} Librería de Node.js que proporciona una API de alto nivel para controlar navegadores Chromium headless para renderizado e impresión PDF.
    \item \textbf{CSS Paged Media (@page):} Módulo del W3C que permite definir propiedades de impresión (márgenes, encabezados de página repetitivos, números de página) para medios paginados.
    \item \textbf{OpenXML (OOXML):} Estándar ISO/IEC 29500 para documentos ofimáticos editables en formato \texttt{.docx}.
\end{itemize}

\section{Responsabilidades}
\begin{itemize}[leftmargin=1.5em, itemsep=2pt]
    \item \textbf{Steeven Toala (Responsable EDT 3.15):} Implementación del motor de renderizado con Puppeteer y generación de Word.
    \item \textbf{Heidi Villavicencio (UI / Calidad):} Maquetación CSS de las plantillas de impresión y verificación frente a plantillas originales.
    \item \textbf{Ing. Paulo Torres (Docente Evaluador):} Aprobación de la fidelidad del documento impreso frente a la norma SGC.
\end{itemize}

\section{Normativa Legal y Técnica}
\subsection{Marco Legal Ecuatoriano}
\begin{itemize}[leftmargin=1.5em, itemsep=2pt]
    \item \textbf{Ley Orgánica de Educación Superior (LOES), Art. 122 y 149.}
    \item \textbf{Normativa Institucional del Sistema de Gestión de Calidad (UTA):} Formatos oficiales \texttt{UTA-SGC-A-2-1-P7-T1} y \texttt{T2}.
\end{itemize}
\subsection{Normas Técnicas Internacionales}
\begin{itemize}[leftmargin=1.5em, itemsep=2pt]
    \item \textbf{ISO 32000-1 (PDF 1.7 Specification):} Estándar internacional para documentos digitales portables.
    \item \textbf{ISO/IEC 29500 (Office Open XML File Formats):} Estándar para formatos de documentos Word editables.
\end{itemize}

\section{Método, Políticas y Contenidos}
\subsection{Enfoque Metodológico}
Se estructuró el servicio \texttt{export.service.js}. Para la salida PDF, el servicio inyecta los datos del documento en una plantilla HTML semántica optimizada con CSS para impresión: márgenes exactos de $2.5\text{ cm}$, repetición automática de cabeceras de tabla en quiebres de página (\texttt{thead \{ display: table-header-group; \}}), numeración de páginas (\textit{``Página X de Y''}) y renderizado headless en Chromium con resolución de 300 DPI. Para la salida Word, se utiliza la librería \texttt{docx} en Node.js que construye programáticamente los párrafos, tablas con bordes y celdas combinadas.

\subsection{Políticas de Calidad de Exportación Documental}
\begin{enumerate}[leftmargin=1.5em, itemsep=2pt]
    \item \textbf{Política de Fidelidad Milimétrica:} El documento exportado debe coincidir exactamente en márgenes, logotipos y disposición con el formato original del SGC.
    \item \textbf{Política de No Huérfanos de Firmas:} La sección de firmas jamás debe quedar aislada en una página en blanco sin contenido que la acompañe (\texttt{page-break-inside: avoid}).
    \item \textbf{Política de Integridad y Código QR:} Todo PDF legalizado debe incluir un código QR que enlaza a la URL de verificación pública de la institución con el hash de autenticidad.
\end{enumerate}

\begin{figure}[H]
    \centering
    \includegraphics[width=0.88\textwidth,height=6.0cm,keepaspectratio]{images/sige_plan_trabajo_a4.png}
    \caption{Captura Real de la Maquetación A4 Idéntica al Formato Oficial UTA-SGC-A-2-1-P7-T1.}
    \label{fig:export_a4}
\end{figure}
\FloatBarrier

\section{Descripción de actividades}
\subsection{Matriz de Parámetros de Exportación en Formatos PDF y DOCX}
\begin{table}[H]
\centering
\small
\renewcommand{\arraystretch}{1.2}
\begin{tabularx}{\textwidth}{|>{\bfseries}m{3.2cm}|m{2.8cm}|X|c|}
\hline
\rowcolor{tableheader}
\textbf{Elemento de Formato} & \textbf{Estándar SGC} & \textbf{Implementación Técnica en el Motor} & \textbf{Conformidad} \\
\hline
Tamaño de Papel & A4 ($210 \times 297\text{ mm}$) & Configurado en Puppeteer con directiva CSS \texttt{@page \{ size: A4; \}} & 100\% \\
\hline
Márgenes de Página & $2.5\text{ cm}$ uniformes & Definidos mediante CSS print media con control milimétrico. & 100\% \\
\hline
Membrete Superior & Logotipo y códigos & Tabla de encabezado repetida en cada página con \texttt{LastPage}. & 100\% \\
\hline
Tabla de Actividades & Celdas con bordes & Tablas con ancho $100\%$ y formato numérico con dos decimales. & 100\% \\
\hline
Bloque de Firmas & 3 casillas alineadas & Contenedor con regla \texttt{break-inside: avoid} e imágenes de firmas. & 100\% \\
\hline
Código QR Verificador & Módulo de seguridad & QR vectorial de $2.5\text{ cm}$ con URL y hash SHA-256 embebido. & 100\% \\
\hline
\end{tabularx}
\caption{Matriz de parámetros técnicos de exportación documental.}
\label{tab:export_matrix}
\end{table}

\subsection{Matriz de Endpoints para Descarga de Documentos Oficiales}
\begin{table}[H]
\centering
\small
\renewcommand{\arraystretch}{1.2}
\begin{tabularx}{\textwidth}{|c|m{3.5cm}|X|c|}
\hline
\rowcolor{tableheader}
\textbf{Método} & \textbf{Ruta URL} & \textbf{Tipo de Salida / Cabecera Content-Type} & \textbf{Descarga} \\
\hline
\texttt{GET} & \texttt{/api/export/:id/pdf} & \texttt{application/pdf} con nombre \texttt{Plan\_Docente\_\{id\}.pdf} & Directa \\
\hline
\texttt{GET} & \texttt{/api/export/:id/word} & \texttt{application/vnd.openxmlformats...} (\texttt{.docx}) & Directa \\
\hline
\texttt{GET} & \texttt{/api/export/:id/preview} & Vista previa en streaming para visor PDF embebido en el navegador. & Streaming \\
\hline
\end{tabularx}
\caption{Endpoints RESTful para generación y descarga de documentos.}
\end{table}

\section{Firmas de Responsabilidad}

\noindent
\renewcommand{\arraystretch}{1.2}
\begin{tabularx}{\textwidth}{|>{\raggedright\arraybackslash}m{2.9cm}|>{\raggedright\arraybackslash}X|>{\centering\arraybackslash}m{4.6cm}|>{\centering\arraybackslash}m{2.3cm}|}
    \hline
    \rowcolor{tableheader}
    \textbf{ACCIÓN} & \textbf{NOMBRE Y CARGO} & \textbf{FIRMA} & \textbf{FECHA} \\
    \hline
    \textbf{Elaborado por:} & Heidi Villavicencio\newline Analista de Calidad / UI-UX & \parbox[c]{4.4cm}{\centering\vspace{4pt}\includegraphics[height=1.5cm,width=4.0cm,keepaspectratio]{images/firma_heidi.png}\vspace{4pt}} & 09/10/2026 \\
    \hline
    \textbf{Revisado por:} & Steeven Toala\newline Líder de Proyecto / Arquitecto & \parbox[c]{4.4cm}{\centering\vspace{4pt}\includegraphics[height=1.5cm,width=4.0cm,keepaspectratio]{images/firma_steveen.png}\vspace{4pt}} & 09/10/2026 \\
    \hline
    \textbf{Aprobado por:} & Ing. Paulo Torres, Mg.\newline Docente Evaluador -- Gestión Proyectos & \parbox[c]{4.4cm}{\centering\vspace{8pt}\textbf{\textcolor{gray!80}{\scriptsize [ PENDIENTE DE APROBACIÓN ]}}\vspace{8pt}} & \textbf{\textcolor{gray!90}{Pendiente}} \\
    \hline
\end{tabularx}

\vspace{0.4cm}

\section{Control de Cambios}

\noindent
\renewcommand{\arraystretch}{1.2}
\begin{tabularx}{\textwidth}{|>{\hsize=0.45\hsize\centering\arraybackslash}X|>{\hsize=1.55\hsize\raggedright\arraybackslash}X|>{\hsize=1.2\hsize\raggedright\arraybackslash}X|>{\hsize=0.8\hsize\centering\arraybackslash}X|}
    \hline
    \rowcolor{tableheader}
    \textbf{Versión} & \textbf{Descripción del cambio} & \textbf{Responsable del cambio} & \textbf{Fecha de aprobación} \\
    \hline
    1.0 & Emisión inicial del entregable según EDT del Tablero Trello & Equipo de Proyecto SIGE & 09/10/2026 \\
    \hline
\end{tabularx}

\vspace{0.4cm}

\section{Anexos}

\subsection*{Anexo 1: Ficha Técnica de la Tarea EDT 3.15 en el Tablero Trello}
\begin{table}[H]
\centering
\small
\renewcommand{\arraystretch}{1.2}
\begin{tabularx}{\textwidth}{|>{\bfseries}m{3.5cm}|X|}
\hline
\rowcolor{tableheader}
\textbf{Parámetro EDT} & \textbf{Detalle Oficial del Tablero Trello} \\
\hline
Código de Tarea & EDT 3.15 \\
\hline
Nombre de la Tarjeta & Motor de exportación oficial en PDF y Word \\
\hline
Fase Correspondiente & FASE 3. IMPLEMENTACIÓN Y DESARROLLO \\
\hline
Período de Ejecución & 07 de octubre al 09 de octubre de 2026 (3 días laborables) \\
\hline
Horas Asignadas & 24 horas hombre netas \\
\hline
Responsables Técnicos & Steeven Toala (ST), Heidi Villavicencio (HV) \\
\hline
Entregable Oficial & \texttt{DocumentExportService.cs} con salidas idénticas al SGC \\
\hline
Estado en el Tablero & 100\% Completado y Verificado en Sprint 4 \\
\hline
\end{tabularx}
\caption{Ficha técnica de la tarea EDT 3.15 en el Tablero Trello.}
\end{table}

\subsection*{Anexo 2: Matriz de Validación de Fidelidad Visual de Exportaciones}
\begin{table}[H]
\centering
\footnotesize
\renewcommand{\arraystretch}{1.15}
\begin{tabularx}{\textwidth}{|c|m{3.5cm}|X|c|}
\hline
\rowcolor{tableheader}
\textbf{Ítem Auditado} & \textbf{Plantilla SGC Oficial} & \textbf{Salida Generada por SIGE} & \textbf{Resultado} \\
\hline
Código de Formato & \texttt{UTA-SGC-A-2-1-P7-T1} & Exacto en esquina superior derecha & 100\% CONFORME \\
\hline
Tipografía & Times New Roman / Arial & Escala tipográfica exacta y pesos respetados & 100\% CONFORME \\
\hline
Total de Páginas & 5 a 6 páginas promedio & Paginación fluida sin desbordamientos & 100\% CONFORME \\
\hline
Escaneo de Código QR & Lectura en dispositivo móvil & Redirecciona a la página de verificación oficial & 100\% CONFORME \\
\hline
\end{tabularx}
\caption{Auditoría de correspondencia frente a formatos originales de la UTA.}
\end{table}

\subsection*{Anexo 3: Lista de Chequeo y Criterios de Aceptación Institucional (SGC ISO 9001:2015)}
\begin{table}[H]
\centering
\footnotesize
\renewcommand{\arraystretch}{1.15}
\begin{tabularx}{\textwidth}{|m{3.2cm}|X|c|c|}
\hline
\rowcolor{tableheader}
\textbf{Criterio de Calidad SGC} & \textbf{Mecanismo de Verificación Técnica} & \textbf{Conformidad} & \textbf{Responsable} \\
\hline
Fidelidad Documental & Exportación idéntica a plantillas oficiales UTA-SGC P7-T1 y P7-T2 & 100\% CONFORME & Heidi Villavicencio \\
\hline
Generación PDF Vectorial & Motor de renderizado basado en Puppeteer para documentos nítidos & 100\% CONFORME & Steeven Toala \\
\hline
Generación DOCX & Conversión a formato editable Word respetando márgenes y tipografía & 100\% CONFORME & Andrea Vásquez \\
\hline
Código QR Institucional & Estampado de QR de verificación que valida autenticidad del documento & 100\% CONFORME & Nixon Hurtado \\
\hline
\end{tabularx}
\caption{Lista de verificación y control de calidad institucional del entregable 3.15.}
\end{table}


\end{document}
